%!TEX root = /Users/lmn/Library/Mobile Documents/com~apple~CloudDocs/lmn@icloud/matematicas/TFG-main/LaTex/main.tex
\chapter{La conexión umbral}

Antes de comenzar, hemos de aclarar parte de la notación y fijar varias convenciones que utilizaremos a lo largo del trabajo. En lo que sigue
\begin{itemize}
    \item $\field$ será un cuerpo de característica cero cualquiera (podemos suponer que es $\C$) y este será el cuerpo base donde trabajaremos.\footnote{También se han dado generalizaciones de la teoría clásica del cálculo umbral cuando $\field$ no es de característica cero. Un ejemplo de esto es \cite{verdoodt}.}
    \item $\F = \formal$ será el anillo de series formales en la variable formal $T$. Además, supondremos conocidas la teoría básica y los principales resultados que a esta conciernen, aunque damos el apéndice \ref{series formales} donde se puede consultar parte de esta información.
    \item Utilizaremos $\poly$ para referirnos a la $\field$-álgebra de polinomios (con la estructura de álgebra usual) en una variable, es decir, $\poly \equiv \field[x]$.
    \item $\gr(p)$ simbolizará el grado del polinomio $p(x)$ y $p'(x)$ simbolizará la derivada, en el sentido usual, del polinomio $p(x)$. En el caso de las derivadas sucesivas, estas se escribirán como $p^{(k)}(x)$.
    \item $\delta_{n,k}$ será la delta de Kronecker, es decir,
        $\delta_{n,k} (x) = \begin{dcases}
            1 & k=n\\
            0 & k \neq n
        \end{dcases} $
    \item $(x)_n$ simbolizará al factorial descendente. Esto es: $(x)_n = x(x-1) \cdots (x-n+1)$
\end{itemize}

%: EL ÁLGEBRA UMBRAL
\section{El álgebra umbral}
\begin{definicion}
    Dado un espacio vectorial $V$, llamaremos \textbf{funcional lineal} a todo elemento $L \in V^*$, donde $V^*$ es el espacio dual de $V$.
    \begin{align*}
        L &\colon V \longrightarrow \field\\
        &v \mapsto \langle L | v\rangle
    \end{align*}
\end{definicion}

Nos centraremos en el estudio de los funcionales lineales sobre el espacio de polinomios $\poly$; es decir, las aplicaciones $L \colon \poly \longrightarrow \field$ tales que $\langle L | p(x) + \lambda q(x) \rangle = \langle L | p(x) \rangle + \lambda \langle L | q(x) \rangle$, para todo $p(x), q(x) \in \poly$ y $\lambda \in \field$. 

\begin{obs}
Nótese que hemos fijado la notación $\langle L | p(x) \rangle$ para simbolizar la acción de $L$ sobre $p(x)$, típica de la física. Más tarde se aclarará el motivo de esta elección. Por otro lado, debido a ahorro del lenguaje y sin que cause confusión, llamaremos a los funcionales lineales de $\poly$ simplemente \textit{\textbf{funcionales}}.
\end{obs}

El primer resultado importante que usaremos continuamente proviene del álgebra lineal.
\begin{teorema}
    Si $\{p_n(x)\}_{n\geq0}$ es una sucesión de polinomios tal que $\gr(p_n) = n$ para todo $n$, entonces dos funcionales $L$ y $M$ son iguales si y solo si $\langle L | p_n(x) \rangle = \langle  M | p_n(x) \rangle$ para todo $n$. En particular, $p_n(x) = x^n$ cumple el teorema.
\end{teorema}
\begin{demo}
    Basta con ver que los polinomios $p_n(x)$ así definidos forman una base (de Hamel) de $\poly$. Que son generadores está claro por ser cada polinomio $p_n(x)$ de grado $n$, por tanto si $p(x)$ es de grado $n$, existe una constante $\lambda \in \field$ tal que $\gr(p - \lambda p_{n}) < n$, y por inducción se concluye que $p$ está en el subespacio generado por $p_{0}, p_{1}, \dots, p_{n}$. Veamos que cada subconjunto finito de estos son linealmente independientes. Sean pues $p_{n_1}(x), \cdots , p_{n_r}(x)$ un subconjunto finito de $\{p_n\}_{n\geq0}$, que, sin pérdida de generalidad, podemos suponer ordenado con $\gr(p_{n_1}(x)) < \dots < \gr(p_{n_r}(x))$, y $\lambda_{n_1}, \dots , \lambda_{n_r} \in \field$ escalares tales que $\lambda_{n_1}p_{n_1}(x) + \cdots + \lambda_{n_r}p_{n_r}(x) = 0$. Si escribimos cada polinomio $p_{n_l}(x)$ como $p_{n_l}(x) = \sum_{k=0}^{n_l}a_{k,n_l}x^k$, igualando coeficientes y teniendo en cuenta el orden dado y que $a_{n_k,n_k} \neq 0$ para todo $k$, obtenemos las ecuaciones
    \begin{align*}
        \lambda_{n_r}a_{n_r,n_r} = 0 \quad \lambda_{n_{r-1}}a_{n_{r-1},n_{r-1}} + \lambda_{n_r}a_{n_{r-1},n_r} = 0 \quad \dots \quad \sum_{k=0}^r\lambda_{n_k}a_{0,n_k} = 0
    \end{align*}
    de donde vemos que $\lambda_{n_1} = \cdots = \lambda_{n_r} = 0$. Luego, son linealmente independientes.
\end{demo}
\begin{corolario}\label{spanning argument}
    Dada $\{p_n(x)\}_{n\geq0}$ sucesión de polinomios tal que $\gr(p_n) = n$ para todo $n$, un funcional está completamente determinado por las constantes $\langle L | p_n(x) \rangle$. En particular, un funcional $L$ está completamente determinado por las constantes $\langle L | x^n \rangle$ y si es el caso que verifica $\langle L | p_n(x) \rangle = 0$ para todo $n$, entonces $L = 0$.
\end{corolario}
Este teorema nos es de suma utilidad ya que nos permite definir un funcional de una forma tan sencilla como dar tan solo su valor sobre las potencias $x^n$. Lo aprovechamos ahora para obtener uno de los teoremas fundamentales del cálculo umbral. Este nos permitirá interpretar cada serie formal como un funcional concreto.

Consideremos la serie formal $T^k = \sum_{n=0}^{\infty} \delta_{n,k} T^{n} \in \F$. En virtud de \ref{spanning argument}, esta serie nos permite definir el funcional $\omega_{T^k}$ dado por
\begin{equation}\label{omegaT^k}
    \langle \omega_{T^k} | x^n \rangle = n! \delta_{n,k}
\end{equation}
Es más que natural extender esta definición usando la linealidad a una serie formal genérica. Así pues, si $F(T) = \sum_{k=0}^\infty a_k T^k \in \F$ es una serie formal cualquiera, esta definirá el funcional $\omega_F$ que viene dado por
\begin{equation} \label{omegaF}
    \langle \omega_F | x^n \rangle = n! \cdot a_n
\end{equation}
Recíprocamente, si consideramos un funcional $L \in \poly^*$, este define la serie formal
\begin{equation}\label{F_L}
    F_L(T) = \sum_{k=0}^\infty \frac{ \langle L | x^k \rangle}{k!} T^k 
\end{equation}
Esta correspondencia entre funcionales y series formales es la que nos permite formular y demostrar el siguiente teorema, base de todo el cálculo umbral.

\begin{teorema}[isomorfismo umbral] 
    La aplicación $\phi \colon  \poly^* \xrightarrow{\sim} \F$ definida por
    % \begin{align*}
    %     L \mapsto \phi (L) = F_L(T) = \sum_{k=0}^\infty \frac{ \langle L | x^k \rangle}{k!} T^k \quad \text{y} \quad \omega_F = \phi^{-1}(F(T)) \mapsfrom F(T) = \sum_{k=0}^\infty a_k T^k
    % \end{align*}
    \[
        L \mapsto \phi (L) = F_L(T) = \sum_{k=0}^\infty \frac{ \langle L | x^k \rangle}{k!} T^k
    \]
	es un isomorfismo $\field$-lineal, con inverso dado por 
	\[
		\omega_F = \phi^{-1}(F(T)) \mapsfrom F(T) = \sum_{k=0}^\infty a_k T^k.
	\]
	donde $\omega_{F}$ es el funcional determinado por $\langle \omega_F | x^n \rangle = n! \cdot a_n$.
A este morfismo $\phi$ lo llamaremos el \textbf{isomorfismo canónico umbral} (abreviadamente, i.c.u).
\end{teorema}
\begin{demo}
    Veamos primero que es biyectiva, para lo que hemos de comprobar que $\phi^{-1}$ es la del enunciado. Por un lado,  $F(T) \longrightarrow \omega_F \longrightarrow F_{\omega_F}(T)$ y $F_{\omega_F}(T) = F(T)$ puesto que $\langle \omega_F | x^n \rangle = n! \cdot a_n$. Mientras que por otro lado, $L \longrightarrow F_L(T) \longrightarrow \omega_{F_L}$ y $\omega_{F_L} = L$ ya que, por \ref{spanning argument}, $L = \omega_{F_L}$ si y solo si $\langle L | x^n \rangle = \langle \omega_{F_L} | x^n \rangle$, lo cual es cierto. Queda así probado que $\phi \circ \phi^{-1} = \phi^{-1} \circ \phi = \operatorname{id}$. Falta ver la linealidad, pero esta es inmediata pues es consecuencia de la linealidad en el primer factor de $\langle \quad | \quad \rangle$; es decir, de que $\langle L + \lambda M | p(x) \rangle = \langle L | p(x) \rangle + \lambda \langle M | p(x) \rangle $ para todo $\lambda \in \field$.
\end{demo}

De esta manera, obtenemos una nueva forma de interpretar a las series formales: como funcionales sobre el espacio de polinomios en una variable. Ahora bien, si realmente queremos tratar a las series como a los segundos, debemos comprobar que el comportamiento respecto al producto de funcionales no difiere con el del producto formal; es decir, una vez probado que $\phi$ es un isomorfismo lineal entre el espacio de funcionales $\poly^*$ y el espacio de series formales $\F$, hemos de responder a la siguiente pregunta: ¿es, además, un isomorfismo de álgebras? La respuesta es que, con la estructura usual de $\field$-álgebra de $\poly^*$, no tiene por qué serlo. Esto es evidente ya que
\begin{equation*}
    \phi(L \cdot M) = F_{L \cdot M}(T) = \sum_{n\geq0} \frac{ \langle L \cdot M | x^n \rangle}{n!} T^n = \sum_{n\geq0} \frac{ \langle L | x^n \rangle \langle M | x^n \rangle}{n!} T^n
\end{equation*}
y, en lo general, $\sum_{n\geq0} \frac{ \langle L | x^n \rangle \langle M | x^n \rangle}{n!} T^n \neq (\sum_{n\geq0} \frac{ \langle L | x^n \rangle}{n!} T^n) (\sum_{n\geq0} \frac{ \langle M | x^n \rangle}{n!} T^n) = \phi(L) \cdot \phi(M)$

No obstante, también es cierto que este isomorfismo induce una estructura de $\field$-álgebra en $\poly^*$ distinta de la que en un principio teníamos y para la cual $\phi$ sí será un isomorfismo de álgebras por la propia construcción.

\begin{definicion}
    Llamaremos \textbf{álgebra umbral} al espacio vectorial $\poly^*$ dotado de la estructura de $\field$-álgebra inducida por $\phi$, es decir, aquella definida por
    \begin{equation}\label{producto funcionales}
        \langle L \cdot M | x^n \rangle = \sum_{k=0}^n\binom{n}{k}\langle L | x^k \rangle \langle M | x^{n-k} \rangle \quad \forall L, M \in \poly^*
    \end{equation}
\end{definicion}
Es una mera comprobación ver que, en efecto, esto nos da una estructura de $\field$-álgebra y en forma de resumen de lo dicho hasta ahora, llegamos al siguiente resultado.

\begin{corolario}
    El i.c.u es un isomorfismo de álgebras entre el álgebra umbral y el álgebra de series formales.
\end{corolario}

En particular, el álgebra umbral cumple las siguientes dos propiedades.
\begin{proposicion}\label{identidad algebra umbral}
    El funcional $e^{0T} = 1$ es la identidad para el producto del álgebra umbral.
\end{proposicion}
\begin{demo}
    Basta con usar el apartado 1 de \ref{spanning argument}.
\end{demo}

\begin{proposicion}\label{producto multiples funcionales}
    Sean $L_1, \dots , L_r \in \poly^*$ funcionales. Entonces, se tiene que
    \begin{equation*}
        \langle L_1 \cdots L_m | x^n \rangle = \sum_{|i|=n}\binom{n}{i_1, \dots , i_m} \langle L_1| x^{i_1} \rangle \dots  \langle L_m | x^{i_m} \rangle
    \end{equation*}
    donde $|i| = i_1 + \dots +i_m$ y $\binom{n}{i_1, \dots , i_m} = \frac{n!}{k_1! \cdots k_m!}$
\end{proposicion}
\begin{demo}
    Se ve aplicando inducción a la definición del producto \eqref{producto funcionales}
\end{demo}

\begin{obs}
    Este último resultado se puede generalizar aún más, pues es válido para todas las sucesiones de polinomios que son de \textit{\textbf{tipo binomial}}. Más tarde las definiremos correctamente y veremos que $p_n(x) = x^n$ es una de estas sucesiones. No obstante, nos es suficiente con conocerlo para el caso particular de $x^n$, tal y como se ha formulado en \ref{producto multiples funcionales}.
\end{obs}

Además de esta particular estructura algebraica, al espacio de funcionales también se le puede dotar de una estructura de espacio topológico con la que el i.c.u será un homeomorfismo entre este espacio y el álgebra de series formales (con la topología formal).
\begin{definicion}\label{topologia de estabilizacion}
Dados dos espacios vectoriales $V$ y $W$, se define la \textbf{topología de estabilización de} $\Hom(V,W)$ tomando sobre $W$ la topología discreta y sobre $\Hom(V,W)$ la topología heredada del producto $W^{V}$. Una base de entornos para $f \in \Hom(V,W)$ consta de los conjuntos
\[
U_A(f) = \{ g \in \Hom(V, W) : f(v) = g(v) \ \forall v \in A\}
\]
donde $A \subseteq V$ recorre los subconjuntos finitos de $V$. Al tratarse de morfismos, es inmediato comprobar que esta topología también está determinada por la base de entornos
\[
U_E(f) = \{ g \in \operatorname{Hom}(V, W) \mid g|_E = f|_E \} = \{ g \in \operatorname{Hom}(V, W) \mid E \subseteq \ker(g - f) \}
\]
donde $E$ recorre los \emph{subespacios} de dimensión finita de $V$.


	
\end{definicion}

Si $V$ tiene una base de Hamel numerable $\{e_{n}\}$, la topología de estabilización en $\Hom(V,W)$ está generada por los entornos $U_{E_{n}}(f)$, donde $E_{n}$ es el subespacio generado por $e_{1}, \dots, e_{n}$. En particular, satisface el primer axioma de numerabilidad y por tanto está determinada por la siguiente condición de convergencia: una sucesión de homomorfismos $\{f_n\}_n \subset \Hom(V,W)$ converge a un homomorfismo $f$ si, dado un elemento $v \in V$, existe un número natural $n_0$ tal que si $n \geq n_0$ entonces $f_n(v) = f(v)$, es decir, cuando $\{f_n(v)\}_{n}$ converge en la topología discreta de $W$, es decir, estabiliza.

En particular, la topología de estabilización en $\poly^{*}$ está generada por los entornos
\[
	U_{n}(L) = \{M \in \poly^{*} : \langle M \mid x^{k} \rangle = \langle L \mid x^{k} \rangle \text{ para } 0 \leq k \leq n \}.
\]




\begin{proposicion}\label{icu homeomorfismo}
    Con la topología de estabilización y la estructura de álgebra umbral, el espacio de funcionales $\poly^*$ es un álgebra topológica completa.
\end{proposicion}

\begin{demo}
    Para ver que es un álgebra topológica basta con probar que tanto la suma como el producto de operadores son aplicaciones continuas, ya que, por la definición del producto de operadores, esto implicaría también que el producto por escalares es continuo. 
    
    Consideramos, entonces, sucesiones de operadores $L_n$ y $M_n$ que converjan a ciertos operadores $L$ y $M$, respectivamente; es decir, para todo natural $m$, existe $n_0$ tal que si $n \geq n_0$, entonces $\langle L_n | x^m \rangle = \langle L | x^m \rangle$ y $\langle M_n | x^m \rangle = \langle M | x^m \rangle$. Por tanto, tendremos que
    \begin{align*}
        \langle L_n + M_n | x^m \rangle & = \langle L_n | x^m \rangle + \langle M_n | x^m \rangle  = \langle L | x^m \rangle + \langle M | x^m \rangle = \langle L + M | x^m \rangle  \\
        \langle L_n \cdot M_n | x^m \rangle & = \sum_{k=0}^m \binom{m}{k}\langle L_n | x^k \rangle \langle M_n | x^{m-k} \rangle = \sum_{k=0}^m \binom{m}{k}\langle L | x^k \rangle \langle M | x^{m-k} \rangle = \langle L \cdot M | x^m \rangle
    \end{align*}
    para un $n$ suficientemente grande y para todo $m \geq 0$. Generalizando esto por linealidad a un polinomio $p(x) \in \poly$ cualquiera, obtenemos que $L_n + M_n$ converge a $L + M$ y $L_n \cdot M_n$ converge a $L \cdot M$, con lo que se concluye.

    Queda probar que es completa. Sea $\{L_n\}_{n}$ una sucesión de Cauchy en $\poly^*$; es decir, dado un polinomio $p(x) \in \poly^*$ existe un $n_0$ tal que si $n, m \geq n_0$ entonces $\langle L_n | p(x) \rangle = \langle L_m | p(x) \rangle $. En particular, para el operador $L$ definido por $\langle L | p(x) \rangle \equiv \langle L_{n_0} | p(x) \rangle $, tenemos que $\langle L_n | p(x) \rangle = \langle L_{n_0} | p(x) \rangle = \langle L | p(x) \rangle$. Luego, $L_n \longrightarrow L$ y se concluye.
\end{demo}

\begin{proposicion}
    Con la topología de estabilización de $\poly^*$ y la topología formal, el i.c.u es un homeomorfismo.
\end{proposicion}

\begin{demo} Hemos de ver que $\phi$ y $\phi^{-1}$ son ambas continuas.
    \begin{itemize}
        \item \textit{$\phi$ es continua:} Sea $\{L_n\}_{n}$ sucesión de funcionales tal que converja a un funcional $L$. Esto es, dado $p(x) \in \poly$ existe un $n_0$ tal que si $n \geq n_0$ entonces $\langle L_n | p(x) \rangle = \langle L | p(x) \rangle$. Así pues, de la definición de $F_{L_n}$ y $F_{L}$, vemos que sus términos $k-$ésimos coinciden para un $n$ lo suficientemente grande y para todo $k \geq 0$. Luego, $F_{L_n}$ converge a $F_L$ en la topología formal y, en consecuencia, $\phi$ es continua.
        \item \textit{$\phi^{-1}$ es continua:} De forma similar, sea $\{F_n(T)\}_{n}$ sucesión de series formales tal que converja a una serie formal $F(T)$, siendo estas series las dadas por $F_n(T) = \sum_{k\geq 0} \frac{a_{n, k}}{k!}T^k$ y $F(T) = \sum_{k\geq 0} \frac{a_{k}}{k!}T^k $. Fijado $k$, o en general, un número finito de $k$, existe un natural $n_0$ tal que $a_{n,k} = a_k$ para estos $k$ siempre que $n \geq n_0$. Luego, por la definición de la aplicación $\phi^{-1}$, obtenemos que
    \begin{equation*}
        \langle \phi^{-1}(F_n(T)) | x^k \rangle = \langle \omega_{F_n} | x^k \rangle = k! \cdot a_{n,k} = k! \cdot a_k = \langle \phi^{-1}(F(T)) | x^k \rangle
    \end{equation*}
    para $n \geq n_0$ y los $k$ elegidos. Luego, por la linealidad, se ve que, para un $n$ suficientemente grande, $\phi^{-1}(F_n(T))$ y $\phi^{-1}(F(T))$ coinciden al aplicarla sobre cualquier polinomio $p(x)$. Luego, $\phi^{-1}(F_n(T)) \longrightarrow \phi^{-1}(F(T))$ y, por tanto, $\phi^{-1}$ es continua. \qedhere 
    \end{itemize} 
\end{demo}

\begin{teorema}\label{topologia formal equivalente a topologia de estabilizacion}
    La topología de estabilización de $\poly^*$ se corresponde con la topología inicial del i.c.u.
\end{teorema}


\begin{demo}
Para $F(T) = \sum_{k=0}^{\infty} a_{k} T^{k} \in \F$, la base de entornos habitual de $F(T)$ en la topología formal consta de los 
\(
	V_{n}(F) = \{G = \sum_{k=0}^{\infty} b_{k}T^{k} \in \F : b_{k} = a_{k} \text{ para } 0 \leq k \leq n\} = F(T) + (T^{n+1}).
\)
Hay que comprobar que los $\phi^{-1}(V_{n}(F))$ generan la topología de estabilización. Como $\phi$ es un homeomorfismo, estos son los mismos que $\phi^{-1}(V_{n}(\phi(L)))$ donde $L$ recorre $\poly^{*}$. Ahora es inmediato comprobar que por definición $\phi^{-1}(V_{n}(\phi(L))) = U_{n}(L) = \{M \in \poly^{*} : \langle M \mid x^{k} \rangle = \langle L \mid x^{k} \rangle \text{ para } 0 \leq k \leq n \}$, que como hemos comentado anteriormente, son generadores de la topología de estabilización.
\end{demo}



Por otro lado, como corolario del teorema \ref{icu homeomorfismo}, conviene tener presente lo siguiente.

\begin{corolario}
    Sean $F(T), G(T) \in \F$ dos series formales. Entonces, $F(T) = G(T) $ si y solo si $\omega_F = \omega_G$.
\end{corolario}

Esto nos lleva a la conclusión de que podemos pensar cada serie formal $F(T)$ como el funcional $\omega_F$ del álgebra umbral. Es por ello que escribiremos indistintamente $\F$ o $\poly^*$ para referirnos a ambos conjuntos. Es más, haciendo abuso de la notación escribiremos $F(T)$ para denotar tanto a la serie formal como al funcional asociado $\omega_F$; cómo esta se comporte dependerá del contexto en el que aparezca y no debe generar confusión alguna. En cualquier caso, trataremos de referirnos a los conjuntos en función de como estemos interpretando las series formales en ese momento. Para cerrar la sección, damos algunos ejemplos básicos de funcionales que nos serán de enorme utilidad. 
\begin{ej}
\begin{enumerate}
    \item El \textbf{funcional derivada $k$-ésima} $T^k$ viene dado directamente por \ref{omegaT^k}
    \begin{equation*}
        \langle T^k| x^n \rangle = n!\delta_{n,k} = \begin{dcases}
            n! & k=n\\
            0 & k \neq n
        \end{dcases}
        \quad =  (x^n)^{(k)}(0)
    \end{equation*}
    Y si tomamos un polinomio $p(x)$ cualquiera
    \begin{equation}\label{funcional derivada}
        \langle T^k| p(x) \rangle = p^{(k)}(0)
    \end{equation}
    \item El \textbf{funcional de evaluación} $e^{\lambda T}$ está dado por
    \begin{equation*}
       \langle e^{ \lambda T} | x^n \rangle = \left\langle \sum_{k = 0}^\infty \frac{\lambda^k}{k!}T^k | x^n \right\rangle = \lambda^n 
    \end{equation*}
    Con lo que, de forma general
    \begin{equation}\label{funcional evaluacion}
        \langle e^{\lambda T} | p(x) \rangle = p(\lambda)
    \end{equation}
    Siguiendo la notación formal, simbolizaremos a esta serie por $E_\lambda (T) \equiv e^{\lambda T}$. En caso de que $\lambda = 0$ se dice que $E_0$ es el \textbf{funcional aumento}. Así pues, la proposición \ref{identidad algebra umbral} se puede reformular diciendo que el funcional aumento es la identidad para el producto del álgebra umbral. Para $\lambda = 1$ escribiremos simplemente $E$.
    \item El \textbf{funcional diferencia} $e^{ \lambda T} -1$ está dado por
    \begin{equation}\label{funcional diferencia}
        \langle e^{\lambda T}-1 | p(x) \rangle = p(\lambda)-p(0)
    \end{equation}
    Simbolizaremos a esta serie por $\Delta_\lambda \equiv e^{\lambda T} - 1$. En particular, para $\lambda = 1$ escribiremos simplemente $\Delta$.
    \item El \textbf{funcional de Abel} $T E_\lambda$ viene dado por
    \begin{equation*}
        \langle T E_\lambda | x^n \rangle = \langle Te^{\lambda T} | x^n \rangle = \left\langle \sum_{k\geq0} \frac{\lambda^k}{k!}T^{k+1} | x^n \right\rangle = n\lambda^{n-1}
    \end{equation*}
    Con lo que
    \begin{equation}\label{funcional de abel}
        \langle Te^{\lambda T} | p(x) \rangle = p'(\lambda)
    \end{equation}
\end{enumerate}
\end{ej}


%%%%%%%%%%%%%%%%%
%%%%%%%%%%%%%%%%%
%%%%%%%%%%%%%%%%%
%%%%%%%%%%%%%%%%%
%%%%%%%%%%%%%%%%%

%: PROPIEDADES DE LOS FUNCIONALES
\section{Propiedades de los funcionales}
Mostramos aquí las propiedades fundamentales de los funcionales lineales sobre $\poly$, consecuencias todas ellas, más o menos directas, de su interpretación como series formales mediante el i.c.u., así como de la notación establecida.

A partir de las definiciones dadas y de las igualdades \ref{funcional derivada} y \ref{omegaF} se obtiene un análogo del teorema de Taylor clásico para las series formales (o funcionales), así como una reescritura del mismo en términos de los funcionales $T^n$.

\begin{proposicion}\label{taylor analogo}
    Sean $F(T) \in \F$ y $p(x) \in \poly$. Se tiene que
    \begin{enumerate}
        \item $\begin{aligned}[t]
            F(T) = \sum_{k \geq 0} \frac{ \langle F(T) | x^k \rangle}{k!}T^k \quad \textnormal{\textbf{(desarrollo de Taylor para funcionales)}}
        \end{aligned}$ 
        \item $\begin{aligned}
            p(x) = \sum_{k \geq 0} \frac{ \langle T^k| p(x) \rangle}{k!}x^k \quad \textnormal{\textbf{(desarrollo de Taylor clásico)}}
        \end{aligned}$
    \end{enumerate}
\end{proposicion}

\begin{proposicion} \label{orden mayor que grado}
    Sean $F(T) \in \F$ y $p(x) \in \poly$. Si $\ord(F) > \gr(p)$, entonces se verifica que $\langle F(T)| p(x) \rangle = 0$
\end{proposicion}
\begin{demo}
    Esto es consecuencia de la definición \ref{omegaF}, pues si $F(T) = \sum_{k = 0}^\infty a_k T^k$ y $p(x) = \sum_{n=0}^m b_n x^n$, siendo $m = \gr(p)$, entonces  $\langle F(T)| p(x)\rangle = \sum_{n = 0}^m n! \cdot a_n \cdot b_n$ y como $\ord(F) > \gr(p)$, $a_n = 0$ para $n \leq m$.
\end{demo}

\begin{proposicion}\label{convergencia intercambio}
    Sea $\{ F_k(T) \}_{k\geq0}$ una sucesión de series formales tal que $\ord(F_k) \rightarrow \infty$ cuando $k\rightarrow \infty$. Entonces, $\sum_{k}a_kF_k(T)$ converge en la topología formal y además se verifica que
    \begin{equation*}
        \left\langle \sum_{k=0}^\infty a_k F_k(T) \biggm| p(x)\right\rangle = \sum_{k=0}^\infty a_k \left\langle F_k(T)| p(x) \right\rangle
    \end{equation*}
    para todo $p(x) \in \poly$. En particular, se cumple cuando $\ord(F_k) = k$
\end{proposicion}
\begin{demo}
Sea $m$ tal que $\ord(F_k) > \gr(p)$ para todo $k > m$. Por la proposición \ref{orden mayor que grado} tenemos que
    \begin{align*}
        \left\langle \sum_{k=0}^\infty a_kF_k(T) \biggm| p(x) \right\rangle & = \left\langle \sum_{k \leq m} a_kF_k(T) + \sum_{k > m} a_kF_k(T) \biggm| p(x) \right\rangle \\& = \left\langle \sum_{k \leq m} a_kF_k(T) \biggm| p(x) \right\rangle = \sum_{k \leq m} a_k\langle F_k(T)| p(x)\rangle = \sum_{k=0}^\infty a_k\langle F_k(T)| p(x)\rangle
    \end{align*}
    donde la última suma es finita ya que $\langle F_k(T) | p(x) \rangle = 0$ para todo $k > m$.
\end{demo}

\begin{proposicion}\label{producto escalares}
    Sean $F(T) \in \F$, $p(x) \in \poly^*$ y $\lambda \in \field$ cualesquiera. Se verifica que
    \begin{equation*}
        \langle F(T)| p(\lambda x)\rangle = \langle F(\lambda T) | p(x)\rangle
    \end{equation*}
\end{proposicion}
\begin{demo}
    Para $(\lambda x)^k$ vemos que $\langle T^n| (\lambda x)^k \rangle = \lambda^k n! \delta_{n,k} = \lambda^n n! \delta_{n,k} = \langle (\lambda T)^n| x^k\rangle$, con lo que el resultado es cierto. Para un polinomio $p(x) = \sum_{k=0}^m b_k x^k$, siendo $m = \gr(p)$, y una serie $F(T) = \sum_{k = 0}^\infty a_k T^k$ cualesquiera, usando las dos proposiciones anteriores y la linealidad vemos el resultado general
    \begin{align*}
        \langle F(T) | p(\lambda x) \rangle & = \left\langle \sum_{k=0}^\infty a_k T^k \biggm | \sum_{n=0}^{m} b_n (\lambda x)^n \right\rangle = \sum_{k=0}^m a_k\left(\sum_{n=0}^m b_n\langle T^n| (\lambda x)^k\rangle \right) \\
        & = \sum_{k=0}^m a_k\left(\sum_{n=0}^m b_n\langle T^n| (\lambda x)^k\rangle \right) = \langle F(\lambda T) | p(x)\rangle \qedhere
    \end{align*}
\end{demo}

\begin{proposicion}\label{inyectividad sucesion F_k}
    Sea $\{ F_k(T) \}_{k\geq0}$ sucesión de series formales tal que $\ord(F_k) = k$. En ese caso, la aplicación $\langle F(T) \rangle \colon \poly \longrightarrow \field^\infty$ dada por $p(x) \mapsto \{ \langle F_k(T) | p(x) \rangle \}_{k \geq 0}$ es inyectiva.
\end{proposicion}
\begin{demo}
    Ya que $\ord(F_k) = k$, sabemos que las series formales $F_k(T)$ forman una pseudobase\footnote{Es decir, una base topológica, donde cada serie se escribe de manera única como una combinación lineal infinita convergente de elementos de la base.} de $\F$. En particular, existen $a_{n,k} \in \field$ tales que $T^n = \sum_{k} a_{n,k}F_k(T)$. Por tanto, por la proposición \ref{convergencia intercambio}
    \begin{equation*}
        p^{(n)}(0) = \langle T^n | p(x) \rangle = \left\langle \sum_{k=0}^\infty a_{n,k}F_k(T) \biggm| p(x) \right\rangle = \sum_{k=0}^\infty a_{n,k} \left\langle F_k(T) | p(x) \right\rangle
    \end{equation*}
    Con lo que si $ \langle F(T) \rangle(p(x)) =\{  \langle F_k(T) | p(x) \rangle \}_{k \geq 0} = \{ 0 \}_{k \geq 0}$, entonces se ha de verificar que $\sum_{k} a_{n,k} \langle F_k(T) | p(x) \rangle = p^{(n)}(0) = 0$. Y ya que esto vale para todo $n$, $p(x) = 0$.
\end{demo}

\begin{proposicion}
    Sea $\{p_k(x)\}_{k \geq 0}$ familia de polinomios que sea un sistema de generadores algebraico para $\poly$. En ese caso, la aplicación $ \langle p(x) \rangle \colon  \poly^* \longrightarrow \field^\infty$ dada por $F(T) \mapsto \{  \langle F(T) | p_k(x) \rangle \}_{k \geq 0}$
    es inyectiva. En particular, el resultado se verifica para los $p_k(x)$ tales que $\gr(p_k) = k$. 
\end{proposicion}
\begin{demo}
    Como la anterior, escribiendo $x^n$ como combinación lineal de los $p_k(x)$. 
\end{demo}

\begin{proposicion}\label{adjuncion M}
    Sean $F(T) \in \F$ y $p(x) \in \poly$ cualesquiera. Entonces, se verifica que
    \begin{equation*}
        \langle \partial F(T) | p(x) \rangle = \langle F(T) | x \cdot p(x)\rangle
    \end{equation*}
\end{proposicion}

\begin{demo}
    Sea $F(T) = \sum_{k=0}^\infty a_k T^k$. Por la linealidad, basta verlo para $p(x) = x^n$. En ese caso 
        \begin{align*}
            \langle \partial F(T)| x^n\rangle & = \sum_{k=1}^\infty ka_k\langle T^{k-1}| x^n\rangle = \sum_{k=0}^\infty (k+1)a_{k+1}\langle T^{k}| x^n\rangle \\
            & = \sum_{k=0}^\infty (k+1)a_{k+1}k! \delta_{n,k} = (n+1)!a_{n+1} = \langle F(T)| x^{n+1}\rangle = \langle F(T)| x\cdot x^n\rangle \qedhere
        \end{align*}
\end{demo}
    

Recordamos que el operador de integración formal $\int$ está determinado por la relación $\partial \int F = F$ y $(\int F)(0) = 0$; en términos de los coeficientes, $\int \sum_{k=0}^{\infty} a_{k} T^{k} = \sum_{k=0}^{\infty} a_{k} \frac{T^{k+1}}{k+1}$.
	
	
\begin{proposicion}\label{adjuncion operador integral}
    Si $F(T) \in \F$ es una serie formal tal que $\ord(F) \geq 1$, entonces
    \begin{equation*}
        \left\langle \frac{F(T)}{T} \biggm| p(x) \right\rangle = \left\langle F(T) \biggm| \int p(x) \right\rangle
    \end{equation*}
\end{proposicion}

\begin{demo}
    Sea $F(T) = \sum_{k=1}^\infty a_kT^k$. Tenemos que
    \begin{align*}
           \left\langle \frac{F(T)}{T} \biggm| x^n \right\rangle
		 = \left\langle \sum_{k=1}^\infty a_kT^{k-1} \biggm| x^n \right\rangle
		&= \sum_{k=0}^\infty a_{k+1} \langle T^{k}| x^n \rangle 
		\\
		&= n! a_{n+1} 
		 = (n+1)!\frac{a_{n+1}}{(n+1)} 
		 = \left\langle F(T) \biggm| \int x^n \right\rangle
    \end{align*}
    Extendiéndolo a todo $p(x) \in \poly$ por linealidad se concluye.
\end{demo}
    
\begin{ej}
    El \textbf{funcional de Bernoulli} $\frac{e^{\lambda T} -1}{T} = \frac{\Delta_\lambda}{T}$ (bien definido, ya que $\ord(e^{\lambda T} - 1) = 1$) será aquel con
    \begin{equation*}
        \left\langle \frac{\Delta_{\lambda}}{T} \biggm| p(x) \right\rangle = \left\langle \Delta_\lambda \biggm| \int p(x) \right\rangle = \left( \int p(x) \right)(\lambda) - \left( \int p(x) \right)(0) = \left( \int p(x) \right)(\lambda)
    \end{equation*}
    Que en el caso de que $\field = \C, \R$ se reduce a
    \begin{equation} \label{funcional Bernoulli}
        \left\langle \frac{e^{\lambda T} -1}{T} \biggm| p(x) \right\rangle = \left\langle \frac{\Delta_\lambda}{T} \biggm| p(x) \right\rangle = \int_{0}^{\lambda} p(u) \, du
    \end{equation}
\end{ej}

Dos familias de funcionales son de especial importancia: estos son los \textit{funcionales delta} y los \textit{funcionales invertibles}. Los utilizaremos cuando hablemos de las \textit{sucesiones de Sheffer}, puesto que los necesitaremos para construirlas.

\begin{definicion}
    Dada $F(T) \in \F$, diremos que $F(T)$ es una \textbf{serie delta} (o \textbf{$\delta$-serie}) si $F(T) \in \maximal$ y $F(T) \not \in \maximal^2$; es decir, cuando $F(T)$ es un invertible respecto de la composición. Escribimos $\F^\delta$ para simbolizar el conjunto de series delta.
\end{definicion}

\begin{definicion}
    Dada un funcional $L \in \poly^*$, diremos que $L$ es un \textbf{funcional delta} (o \textbf{$\delta$-funcional}) si al pensarlo como una serie (es decir, $\phi(L)$) es una serie delta. Análogamente, diremos que $L$ es un \textbf{funcional invertible} si al pensar $L$ como una serie formal esta es una serie invertible.
\end{definicion}

\begin{proposicion}
    Sea $F(T) \in \F$. Entonces
    \begin{enumerate}
        \item $F(T) \in \F^{\delta} \iff \langle F(T)| 1\rangle = 0, \langle F(T)| x \rangle \neq 0$
        \item $F(T) \in \F^* \iff \langle F(T)|1 \rangle \neq 0$
    \end{enumerate}
\end{proposicion}

\begin{demo}
    Se sigue de las definiciones y de las propiedades para series formales.
\end{demo}

%%%%%%%%%%%%%%%%%%%%%
%%%%%%%%%%%%%%%%%%%%%
%%%%%%%%%%%%%%%%%%%%%
%%%%%%%%%%%%%%%%%%%%%
%%%%%%%%%%%%%%%%%%%%%

%: OPERADORES LINEALES
\section{La relación con los operadores lineales}

\begin{definicion}
    Sea $V$ un espacio vectorial cualquiera. Llamamos \textbf{operador lineal} de $V$ a todo elemento $\Lambda$ de $\End(V)$.
    \begin{align*}
        \Lambda & \colon V \longrightarrow V\\
        & v \mapsto \Lambda v
    \end{align*}
\end{definicion}

\begin{obs}
    De forma similar a lo que hacíamos con los funcionales, a los operadores lineales de $\End(\poly)$ los llamaremos simplemente \textbf{operadores}. Por otro lado, la convención seguida para simbolizar la acción de un operador $\Lambda$ sobre un polinomio $p(x)$ es la de escribir la yuxtaposición del símbolo del operador con el del polinomio: $\Lambda p(x)$. Esto no debe causar confusión con el producto de operadores y, además, es el motivo principal por el que escogimos la notación $\langle F(T)| p(x)\rangle$ para la acción de los funcionales.
\end{obs}

Previamente hemos identificado el espacio de series formales, $\F$, con el espacio de funcionales, $\poly^*$, y como consecuencia hemos dotado a $\poly^*$ de una estructura de álgebra, inducida por el álgebra de las series formales. Una nueva interpretación de las series formales nos permite ahora dotar a $\poly$ de una muy útil estructura de $\F$-módulo.

Para ello, definimos una acción de $\F$ sobre $\poly$ declarando que $T$ actúa mediante la derivación habitual $\partial_{x}$ en polinomios, determinada por
% 
% \begin{align} 
%     & \F \longrightarrow \End(\poly) \nonumber \\
%     & T^k \longleftrightarrow \partial_{x}^k
% \end{align}
% donde es importante recordar que $\partial_{x}^k$ actúa sobre $\poly$ como
\begin{align}
\label{T como derivada}
	\partial_{x}^k x^n = 
	\begin{dcases}
        (n)_kx^{n-k} & k \leq n
		\\
        0 & k > n
    \end{dcases}
\end{align}
Extendemos por linealidad mediante la aplicación $\varphi$ definida por
\begin{align*}
    \varphi \colon & \F \longrightarrow \End(\poly)\\
    & \sum_{k = 0}^\infty \frac{a_k}{k!}T^k \mapsto \sum_{k = 0}^\infty \frac{a_k}{k!}\partial_{x}^k
\end{align*}
que para cada $F(T) \in \F$ asigna el operador $F(\partial_{x})$ determinado por
\begin{align}
\label{accion F(d)xn}
    F(\partial_{x})x^n =\sum_{k = 0}^n \binom{n}{k}a_k x^{n-k}
\end{align}
Es importante destacar que hemos escrito los elementos de $\F$ en notación exponencial (es decir, con un factor $\frac{1}{k!}$). Esto es así puesto que nos resultará más conveniente lidiar con los coeficientes binomiales que con los factoriales descendentes y en lo que sigue esta será la convención que utilizaremos.\footnote{Esta constante se puede cambiar por otra y obtener una teoría similar. Es de especial importancia el cálculo q-umbral. Véase \cite{roman3} y \cite{roman4}.}


Debemos verificar que efectivamente tenemos una acción que dota a $\poly$ de estructura de $\F$-módulo. Esto se puede ver de varias maneras equivalentes, que destacamos por separado.


\begin{proposicion}
    La aplicación $\varphi \colon \F \longrightarrow \End(\poly)$ anterior es un morfismo de álgebras; es decir, para toda $F(T), G(T) \in \F$ se cumple que
    \begin{equation*}
        \varphi(F(T) \cdot G(T)) = \varphi(F(T)) \circ \varphi(G(T))
    \end{equation*}
\end{proposicion}

\begin{demo}
    Que es un morfismo lineal y que $\varphi(1) = 1$ es trivial por su definición. Así mismo, para ver que respeta el producto basta con probarlo para las potencias $T^n$, es decir, basta con ver
    \begin{equation*}
        \varphi(T^{i} \cdot T^{j}) = \varphi(T^{i}) \circ \varphi(T^{j}) \iff \partial_x^{i+j} = \partial_x^{i} \circ \partial_x^{j}
    \end{equation*}
    lo cual es cierto evidentemente.
\end{demo}

De forma similar a lo que hacíamos con los funcionales, utilizaremos el símbolo $T$ para hablar tanto de la variable formal como del operador $\partial_{x}$, con lo que escribiremos siempre $F(T)$ en lugar del operador asociado $\varphi(F(T))$. En caso de no indicarlo, simplemente el contexto ha de servir para diferenciar como se está considerando.

En resumen, dada $F(T)$ la escribiremos como $F(T) = \sum_k \frac{a_k}{k!}T^k$ y esta actúa sobre $\poly$ como
\begin{align}\label{serie como operador}
    F(T) \colon & \poly \longrightarrow \poly \nonumber \\
    & x^n \mapsto F(T)x^n = \sum_{k=0}^n \binom{n}{k}a_kx^{n-k}
\end{align}

\begin{obs}
    En particular, con esta notación la proposición anterior se traduce a que, como operadores
    \begin{equation*}
        F(T) \cdot G(T) = F(T) \circ G(T) \quad \forall F(T), G(T) \in \F
    \end{equation*} 
    Así pues, simbolizaremos a cualquiera de estos por la habitual omisión del símbolo correspondiente (\textit{i.e:.} $F(T)G(T)$) y no deberá existir confusión.
\end{obs}


El morfismo $\varphi$ da a $\poly$ estructura de $\F$-módulo, mediante la acción $\varphi(F(T)) p(x)$. Lo expresamos y probamos también de la manera habitual.


\begin{teorema}
    El producto definido por
    \begin{align*}
        \star \colon & \F \times \poly \longrightarrow \poly\\
        & (F(T), p(x)) \mapsto F(T) \star p(x) \equiv F(T) p(x)
    \end{align*}
    siendo $F(T) p(x)$ definido por \eqref{serie como operador}, dota a $\poly$ de una estructura de $\F$-módulo. Además, es un $\F$-módulo fiel.
\end{teorema}
\begin{demo}
    A partir de la definición de $\star$ es claro que
    \begin{align*}
        (F(T) + G(T)) \star p(x) & = F(T) \star p(x) + G(T) \star p(x)\\
        F(T) \star (p(x) + q(x)) & = F(T) \star p(x) + F(T) \star q(x)\\
        1 \star p(x) & = p(x)
    \end{align*}
    Para ver que $F(T) \star (G(T) \star p(x)) = (F(T) \cdot G(T)) \star p(x)$ basta con tener en cuenta que
    \begin{align*}
        T^k T^j \star p(x) = T^{k+j} \star p(x) = \partial^{k+j} \star p(x) = \partial^{k+j} p(x) = \partial^k(\partial^j p(x)) = T^k \star (T^j \star p(x)) 
    \end{align*}
    y por la linealidad se concluye que es, efectivamente, un $\F$-módulo.

    Queda probar que el módulo es fiel, es decir, comprobar que el conjunto de los anuladores $\text{Ann}(\poly) = \{ F(T) \in \F : F(T) \star p(x) = 0, \forall p(x) \in \poly\}$ es nulo. Consideremos, pues, una serie formal $F(T) = \sum_{k=0}^\infty \frac{a_k}{k!}T^k \in \text{Ann}(\poly)$. Aplicándola a cada $x^n$ como operador obtenemos que
    \begin{align*}
        F(T) \star x^0 & = a_0 = 0 \implies a_0 = 0\\
        F(T) \star x^1 & = a_0x + a_1 = 0 \implies a_1 = 0\\
        F(T) \star x^2 & = a_0x^2 + 2a_1x + a_2 = 0 \implies a_2 = 0\\
        &\vdots
    \end{align*}
    de donde vemos que $a_0 = a_1 = \cdots = a_n = \cdots = 0$ y se concluye.
\end{demo}


Los operadores lineales $F(T)$ cumplirán propiedades análogas a la de los funcionales. Listamos algunas de estas, aunque no damos su demostración puesto que son muy similares a las dadas para estos útlimos.

\begin{proposicion}
    Si $\ord(F) > \gr(p)$, entonces $F(T)p(x) = 0$.
\end{proposicion}

\begin{proposicion}
    Si $\{F_k(T)\}_{k \geq 0}$ es una sucesión de series formales verificando que $\ord(F_k) \rightarrow \infty$, entonces
    \begin{equation*}
        \left( \sum_{k=0}^\infty a_kF_k(T) \right)p(x) = \sum_{k=0}^\infty a_k \left( F_k(T)p(x) \right)
    \end{equation*}
    para todo polinomio $p(x) \in \poly$. En particular, un operador está completamente determinado por su acción sobre los elementos de una pseudobase de $\F$.
\end{proposicion}


\begin{proposicion}
    Si $\{ F_k(T) \}_{k \geq 0}$ es un sistema de generadores topológico para $\F$, entonces la aplicación $\{F(T)\} \colon  \poly \longrightarrow \field^\infty$ dada por $p(x) \mapsto \{ F_k(T)p(x)\}_{k \geq 0}$ es inyectiva.
\end{proposicion}

\begin{proposicion}
    Si $\{ p_k(x) \}_{k \geq 0}$ es un sistema de generadores algebraico para $\poly$, entonces la aplicación $\{p(x)\} \colon \poly \longrightarrow \field^\infty$ dada por $p(x) \mapsto \{ F(T)p_k(x)\}_{k \geq 0}$ es inyectiva.
\end{proposicion}

Uno de los resultados claves en lo que a estos operadores respecta es el de la relación de adjunción que guardan con el producto del álgebra formal.
\begin{teorema}[de adjunción]\label{adjuncion serie operador y multiplicacion}
    Sean $F(T), G(T) \in \F$. Para todo $p(x) \in \poly$ se verifica
    \begin{equation*}
        \langle F(T)G(T) | p(x) \rangle = \langle F(T) | G(T)p(x) \rangle = \langle G(T) | F(T)p(x) \rangle
    \end{equation*}
\end{teorema}

\begin{demo}
    Por la linealidad, esto se reduce a ver que $\langle T^k T^j | p(x) \rangle = \langle T^k | T^j p(x) \rangle$. En efecto, debido a la conmutatividad del producto de series formales se verifica que
    \begin{equation*}
        \langle T^k T^j | p(x) \rangle  = \langle T^{k+j} | p(x) \rangle = p^{(k+j)}(0) = (\partial^j p)^{(k)}(0) = (T^j p)^{(k)}(0) = \langle T^k | T^j p(x) \rangle \qedhere
    \end{equation*}
\end{demo}

\begin{corolario}
    Sea $F(T) \in \F$ y $p(x) \in \poly$ cualesquiera. Entonces, se verifica que
        \begin{align*}
            \langle F(T) | p'(x)\rangle = \langle T\cdot F(T)| p(x)\rangle
       \end{align*}
    Equivalentemente, se verifica que $\langle F(T) | p(x)\rangle = \left\langle T\cdot F(T)| \int p(x) \right\rangle$
\end{corolario}

\begin{corolario}
    Dada $F(T) \in \F$ y $p(x) \in \poly$ se tiene que $ \langle  F(T) | p(x) \rangle = (F(T)p(x))(0)$. Es decir, se verifica que el diagrama
    % $\vcenter{\hbox{
\[
	\begin{tikzcd}
    {\poly} & {\poly}\\
    & {\field}
    \arrow["{F(T)}", from=1-1, to=1-2]
    \arrow["{\langle F(T) | \quad \rangle}"', sloped, above, from=1-1, to=2-2]
    \arrow["{E_0}", from=1-2, to=2-2]
    \end{tikzcd}
\]
	% }}$
es conmutativo.
\end{corolario}
Si bien la asignación \eqref{T como derivada} es biunívoca de forma trivial, a diferencia de lo que ocurría con el i.c.u, la aplicación $\varphi$ \underline{\textbf{no}} es una aplicación biyectiva; es decir, existen operadores que no provienen de ninguna serie formal. Esto se debe a que $\gr(\partial_{x} p) \leq \gr(p)$, luego, por linealidad, se tiene que $\gr(F(T)p(x)) \leq \gr(p)$, es decir, que todos los operadores que son series formales actúan bajando o manteniendo el grado.

Así pues, el operador \textit{multiplicación por $x$}, dado por $p(x) \mapsto xp(x)$, en tanto que cumple que $\gr(xp(x)) \geq \gr(p)$, es un ejemplo de un operador que no proviene de ninguna serie formal.

\begin{obs}\label{multiplicacion por x operador}
    El operador \textbf{\textit{multiplicación por $x$}} nos aparecerá posteriormente. En lo que sigue, lo escribiremos abreviadamente como $M$. Hacemos esto con el objetivo de evitar la posible confusión que podría crearse con la multiplicación de series formales y polinomios. De una forma más general, dado un polinomio $p(x)$, simbolizamos el operador \textit{\textbf{multiplicación por $p(x)$}} como $M(p(x))$. En particular, la proposicón \ref{adjuncion M} se reescribe como
    \begin{equation}\label{adjuncion multiplicar x}
        \langle \partial F(T) | p(x) \rangle = \langle F(T) | M p(x)\rangle
    \end{equation}
    de donde surge una nueva aparente relación de adjunción entre los operadores $M$ y $\partial$. Todo esto se esclarecerá en el capítulo tercero. 
\end{obs}

Mediante lo ya dicho, se ha probado la siguiente proposición.

\begin{proposicion}
    La aplicación $\varphi \colon \F \hookrightarrow \End(\poly)$ es un homomorfismo de anillos inyectivo pero no epiyectivo, con imagen
    \begin{equation*}
        \text{Im}(\varphi) = \{ \Lambda \in \End(\poly) : \exists F(T) \in \F \text{ \textnormal{con} } \Lambda p(x) = F(T)p(x), \forall p(x) \in \poly \}
    \end{equation*}
\end{proposicion}

Nos interesa ahora, por lo tanto, dar una caracterización de la imagen de este morfismo.

\begin{definicion}
    Diremos que un operador $L \in \End(\poly)$ tiene \textbf{forma de serie formal} si $L \in \text{Im}(\varphi)$.
\end{definicion}

\begin{definicion}
    Diremos que un operador $F(T) \in \text{Im}(\varphi)$ es un \textbf{operador delta} (o \textbf{$\delta$-operador}) si como serie formal $F(T)$ es una serie delta. Análogamente, diremos que $F(T)$ es un \textbf{operador invertible} si como serie formal $F(T)$ es una serie invertible. 
\end{definicion}

\begin{lema}
    Si $F(T)$ es un operador delta y $\Lambda \in \End(\poly)$ conmuta con $F(T)$, es decir, $F(T) \circ \Lambda = \Lambda \circ F(T)$, entonces $\gr(\Lambda p(x)) \leq \gr(p)$ para todo $p(x) \in \poly$
\end{lema}
\begin{demo}
    Como venimos haciendo, por la linealidad basta con probarlo para los polinomios $x^n$. Por contradicción, supongamos que existe un número natural $m$ tal que $\gr(\Lambda x^m) \geq m+1$. Ya que $\ord(F) = 1$, tenemos que $F(T) = \sum_{k = 1}^\infty \frac{a_k}{k!}T^k$, de forma que como operador verificará que $F(T)x^n = \begin{dcases}
        \sum_{k=1}^n \binom{n}{k}a_k x^k = n a_1 x^{n-1} + \cdots + a_n \quad & n \geq 1\\
        0 & n = 0
        \end{dcases}$ 
    
    Por tanto, $\gr(F(T)p(x)) = \gr(p) - 1$ si $\gr(p) \geq 1$. A partir de aquí deducimos que como $\ord(F^{m+1}) = m+1$, entonces $F^{m+1}(T)x^m = 0$. Y en consecuencia se tiene que $F^{m+1}(T) (\Lambda x^m) = \Lambda (F^{m+1}(T)x^m) = 0$. Pero, por otro lado, como hemos supuesto que $\gr(\Lambda x^m) \geq m+1$ tendremos que $\gr(F^{m+1}(T) (\Lambda x^m)) \geq 0$ y por tanto $F^{m+1}(T) (\Lambda x^m) \neq 0$, lo cual es imposible. Con lo que $\gr(\Lambda x^n) \leq \gr(x^n)$ necesariamente y se concluye.
\end{demo}

\begin{teorema}[caracterización Im($\varphi$)]
    Un operador $\Lambda \in \End(\poly)$ tiene forma de serie formal si y solo si conmuta con el operador derivada (i.e: $T \circ \Lambda = \Lambda \circ T$).
\end{teorema}

\begin{demo}
    Es claro que los elementos de $\text{Im}(\varphi)$ conmutan con $T$, puesto que conmutan con todas las series formales, luego en particular lo hacen con $T$. Recíprocamente, si $\Lambda$ conmuta con $T$, entonces veamos que $\Lambda$ tendrá la forma de $F(T)$ para la serie formal $F(T)$ siguiente:
    \begin{equation*}
        F(T) = \sum_{k=0}^\infty \frac{\langle 1| \Lambda x^k\rangle}{k!}T^k
    \end{equation*}
    Por el lema previo tenemos que $\gr(\Lambda x^n) \leq n$. En consecuencia, el desarrollo de Taylor de \ref{taylor analogo} nos da que $\Lambda x^n = \sum_{k=0}^n \frac{\langle T^k| \Lambda x^n\rangle}{k!}x^k$. Utilizando ahora las propiedades demostradas previamente junto a la hipótesis de conmutación con $T$, vemos entonces que $\Lambda x^n$ cumple que $\langle T^k| \Lambda x^n\rangle = \langle 1| (T^k \circ \Lambda) x^n\rangle = \langle 1| (\Lambda \circ T^k) x^n\rangle = \langle 1| \Lambda [(n)_k x^{n-k}]\rangle = (n)_k \langle 1| \Lambda x^{n-k}\rangle$. Por tanto, podemos reescribir a este polinomio como
    \begin{equation*}
        \Lambda x^n = \sum_{k=0}^n \frac{(n)_k}{k!}\langle 1| \Lambda x^{n-k}\rangle x^k = \sum_{k=0}^n \binom{n}{k}\langle 1| \Lambda x^{n-k}\rangle x^k = \sum_{k=0}^n \binom{n}{k}\langle 1| \Lambda x^{k}\rangle x^{n-k} = F(T)x^n 
    \end{equation*}
    donde hemos usado que $\binom{n}{k} = \binom{n}{n-k}$, y llegamos así a que $\Lambda = F(T)$ usando \ref{spanning argument}.
\end{demo}

\begin{corolario}\label{conmutacion G(T)}
    Dado un operador delta $G(T)$, $\Lambda \in \End(\poly)$ tiene la forma de la serie formal $F(T)$ si y solo si conmuta con $G(T)$
\end{corolario}

\begin{demo}
    Si $\Lambda$ tiene forma de serie formal entonces conmuta con toda serie formal, en particular con $G(T)$. Recíprocamente, si $\Lambda \circ G(T) = G(T) \circ \Lambda$, como $G(T)$ es un operador delta, las series formales $G^k(T)$ forman una pseudobase de $\F$. En concreto, existen escalares $a_k \in \F$ tales que $T = \sum_{k = 0}^\infty a_k G^k(T)$. Por lo tanto, la conmutatividad implica que
    \begin{align*}
        \Lambda (T x^n) & = \Lambda\left(\sum_{k=0}^\infty a_k G^k(T) x^n\right) = \Lambda \left(\sum_{k=0}^n a_k G^k(T) x^n\right) = \sum_{k=0}^n a_k (\Lambda \circ G^k(T)) x^n\\
        & = \sum_{k=0}^n a_k (G^k(T) \circ \Lambda) x^n = \left(\sum_{k=0}^\infty a_k G^k(T)\right) (\Lambda x^n) = T (\Lambda x^n)
    \end{align*}  
    donde la segunda igualdad es consecuencia de que $\ord(G^k) = k$ y la penúltima se debe a que $\gr(\Lambda x^n) \leq n$ por el lema previo.
\end{demo}

\begin{corolario}
    $\Lambda \in \End(\poly)$ tiene forma de serie formal si y solo si conmuta con algún operador de traslación $E_{\lambda}(T)$ con $\lambda \neq 0$
\end{corolario}

\begin{demo}
    $\Lambda$ conmuta con $E_{\lambda}(T)$ si y solo si $\Lambda$ conmuta con el operador diferencia $\Delta_{\lambda}$. Como el operador diferencia es un operador delta siempre que $\lambda \neq 0$, por el corolario anterior se concluye.   
\end{demo}

Por otro lado, conviene que dotemos al espacio de operadores de una estructura, además, de espacio topológico, que haga que la aplicación $\varphi$ sea continua, también de una forma similar a como lo hacíamos con los funcionales. Esta topología se corresponde nuevamente con la topología de estabilización. Las demostraciones del siguiente par de enunciados son casi idénticas y quedan omitidas.

\begin{proposicion}
    Con la topología de estabilización y la estructura de álgebra usual, el espacio de operadores $\End(\poly)$ es un álgebra topológica completa.
\end{proposicion}

\begin{proposicion}
    Con la topología formal y la topología de estabilización, la aplicación $\varphi$ que asigna a cada serie formal su operador asociado es continua.
\end{proposicion}


También tenemos el siguiente resultado análogo.

\begin{proposicion}
    % La topología de estabilización en $\End(\poly)$ se corresponde con la topología inicial de la aplicación $\varphi$.
    La topología formal en $\F$ se corresponde con la topología inicial de la aplicación $\varphi$, dotando a $\End(\poly)$ de la topología de estabilización.
\end{proposicion}



\begin{demo}
% Como las topologías correspondientes satisfacen el primer axioma de numerabilidad, podemos trabajar más comodamente con sucesiones.
Dada una sucesión $\{F_{n}(T)\}_{n} \subseteq \F$ y $F(T) \in \F$, $F_{n}(T) \to F(T)$ en la topología inicial de $\varphi$ si y solo si $F_{n}(\partial_{x}) \to F(\partial_{x})$ en la topología de estabilización de $\End(\poly)$, que significa que dado $m \geq 0$ existe $n_{0}$ tal que $F_{n}(\partial_{x}) x^{m} = F(\partial_{x}) x^{m}$ si $n \geq n_{0}$. Si $F_{n}(T) = \sum_{k=0}^{\infty} a_{n,k} T^{k}$ y $F(T) = \sum_{k=0}^{\infty} a_{k} T^{k}$, recordando~\eqref{accion F(d)xn} esta relación equivale a
\[
	\sum_{k=0}^{m} \binom{m}{k} a_{n,k} x^{m-k} = \sum_{k=0}^{m} \binom{m}{k} a_{k} x^{m-k} \text{ si } n \geq n_{0},
\]
y comparando coeficientes concluimos que $a_{n,k} = a_{n}$ para $0 \leq k \leq m$ si $n \geq n_{0}$. Esto es precisamente la convergencia $F_{n}(T) \to F(T)$ en la topología formal.
\end{demo}



















Recuperamos ahora los ejemplos de funcionales dados con anterioridad con el objetivo de dar su intepretación como operadores y cerramos así esta sección, junto al primer capítulo.

\begin{ej}
    \begin{enumerate}
        \item El \textbf{operador derivada $k$-ésima} $T^k$ viene dado por \ref{T como derivada} y para un polinomio genérico $p(x) \in \poly$ será
        \begin{equation} \label{operador derivada}
            T^k p(x) = p^{(k)}(x)
        \end{equation}
        \item El \textbf{operador traslación} $E_\lambda$ cumple que
        \begin{equation*}
            E_{\lambda} x^n = \sum_{k=0}^\infty \frac{\lambda^k}{k!} T^k x^n = \sum_{k=0}^n \binom{n}{k}\lambda^k x^{n-k} = (x+ \lambda )^n
        \end{equation*} 
        Con lo que, para un polinomio genérico $p(x)$
        \begin{equation} \label{operador traslacion}
            E_{\lambda} p(x) = p(x+\lambda)
        \end{equation}
        \item El \textbf{operador diferencia} $\Delta_{\lambda}$ en este caso será
        \begin{equation}\label{operador diferencia}
            \Delta_{\lambda} p(x) = p(x+ \lambda) - p(x)
        \end{equation}
        \item El \textbf{operador de Abel} $TE_{\lambda}$ será, por el producto de operadores
        \begin{equation}\label{operador de Abel}
            TE_{\lambda}p(x) = p'(x+ \lambda)
        \end{equation}
        \item El operador $\frac{\Delta_{\lambda}}{T}$ verifica que
        \begin{align*}
            \frac{\Delta_{\lambda}}{T}x^n & = \left(\sum_{k=0}^\infty \frac{\lambda^{k+1}}{(k+1)!}T^k\right) x^n = \sum_{k=0}^n \binom{n+1}{k+1} \frac{\lambda^{k+1}}{(n+1)}x^{n-k}\\
            & = \frac{1}{n+1}\sum_{k=1}^{n+1} \binom{n+1}{k} \lambda^k x^{n+1-k} = \frac{1}{n+1}((x+ \lambda )^{n+1} - x^{n+1}) = \int\limits_{x}^{x+ \lambda} u^n \, du
        \end{align*}
        Por tanto, para un polinomio genérico $p(x)$
        \begin{equation}\label{operador integral media}
           \frac{\Delta_{\lambda}}{T}p(x) = \int\limits_{x}^{x+ \lambda} p(u) \, du
        \end{equation}
    \end{enumerate}
\end{ej}
%%%%%%%%%%%%%%%%%%%%%
%%%%%%%%%%%%%%%%%%%%%
%%%%%%%%%%%%%%%%%%%%%
%%%%%%%%%%%%%%%%%%%%%
%%%%%%%%%%%%%%%%%%%%%
